\documentclass[11pt]{article}
\usepackage{amsmath,amssymb,amsthm,hyperref}
\newtheorem{theorem}{Theorem}
\newtheorem{lemma}[theorem]{Lemma}
\title{Weighted collision bounds for triple minima and six-window swaps}
\author{Research note; independently reviewed by two other agents}
\date{30 September 2026}
\begin{document}
\maketitle

Let $\rho_1,\ldots,\rho_M$ be permutations of $[N]$, with arbitrary
nonnegative weights $\mu_b$ summing to one. Repetitions are allowed.
Put $\kappa=\sum_b\mu_b^2$ and $m=N-1$. The number $1/\kappa$ is
the effective support size used below; the actual support may be much
larger.

\begin{lemma}[Averaging the centered Gram collision term]\label{gram}
For each pivot $i$, let $x_{b,i}\in\{0,1\}^{m}$ have coordinates
$\mathbf1\{\rho_b(i)<\rho_b(j)\}$ for $j\ne i$, and put
\[
 P=I-m^{-1}{\bf1}{\bf1}^{T},\qquad
 R_i=\sum_b\mu_b(Px_{b,i})(Px_{b,i})^T.
\]
If $N\ge3$, some pivot $i$ satisfies
\[
 \|R_i\|_F\ge\frac m{\sqrt{60}}\sqrt\kappa.
\]
\end{lemma}
\begin{proof}
All the weights are nonnegative, so the Gram expansion gives
\[
 \|R_i\|_F^2
 =\sum_{b,c}\mu_b\mu_c
          \langle Px_{b,i},Px_{c,i}\rangle^2
 \ge\sum_b\mu_b^2\|Px_{b,i}\|^4.
\]
For fixed $b$, as $i$ ranges over the labels, the number $k$ of ones in
$x_{b,i}$ ranges over $0,1,\ldots,m$ exactly once. Its centered squared
norm is $k(m-k)/m$. Therefore
\[
 \frac1N\sum_i\|Px_{b,i}\|^4
 =\frac1{m+1}\sum_{k=0}^{m}\frac{k^2(m-k)^2}{m^2}
 =\frac{(m-1)(m^2+1)}{30m}\ge\frac{m^2}{60}.
\]
Averaging the preceding Gram inequality over pivots proves the lemma.
\end{proof}

\begin{theorem}\label{main}
Define
\[
 \epsilon=\max_{i,j,k\text{ distinct}}
  \left|\Pr(\rho(i)<\rho(j),\rho(i)<\rho(k))-\frac13\right|.
\]
Then
\[
 \epsilon\ge\frac{\sqrt\kappa}{\sqrt{60}}
                   -\frac2{3\sqrt{N-1}}.
\]
For $\delta(a,b,c)=\frac12-\frac32\Pr(\rho(b)<\rho(a),\rho(b)<\rho(c))$,
let $\eta$ be the maximum absolute change in the sum of four consecutive
triple scores under a middle-pair swap in a six-label window. For
$N\ge6$ one has
\[
 \eta\ge\frac3{2\sqrt{60}}\sqrt\kappa
                 -\frac4{\sqrt{N-1}}-\frac{128}{N}.
\]
In particular, if $N\ge C/\kappa$ for a sufficiently large absolute
constant $C$, then some such swap has magnitude at least
$c\sqrt\kappa$.
\end{theorem}
\begin{proof}
The uncentered Gram matrix $G_i=\sum_b\mu_bx_{b,i}x_{b,i}^T$ has
diagonal $p_{ij}=\Pr(\rho(i)<\rho(j))$ and off-diagonal triple-minimum
probabilities. Thus
\[
 G_i=\tfrac13{\bf1}{\bf1}^{T}+D_i+E_i,
 \qquad |(D_i)_{jj}|\le2/3,\quad
 (E_i)_{jj}=0,\quad |(E_i)_{jk}|\le\epsilon.
\]
Multiplication on both sides by $P$ removes the constant matrix, giving
$\|R_i\|_F\le(2/3)\sqrt m+\epsilon m$. Combine with
Lemma~\ref{gram} to obtain the first assertion.

For the second assertion, apply the small-swap coboundary lemma from
the companion note \emph{Deterministic six-window discrepancy for short
permutation families}. It supplies an antisymmetric array $a_{ij}$
with $|a_{ij}|\le3/2$ such that
\[
 |\delta(j,i,k)-a_{ij}-a_{ik}|\le\tau:=\eta+128/N.
\]
Consequently, for every pivot $i$,
\[
 G_i=A_i+D_i+E_i,\qquad
 (A_i)_{jk}=\frac13-\frac23(a_{ij}+a_{ik}),
\]
where the formula for $A_i$ also specifies its diagonal,
\[
 (D_i)_{jj}=p_{ij}-\frac13+\frac43a_{ij},\qquad
 |(D_i)_{jj}|\le\frac83,
\]
and $E_i$ has zero diagonal and off-diagonal entries at most
$2\tau/3$ in absolute value. Every summand of the additive matrix
$A_i$ has a constant vector on one side, so $PA_iP=0$. It follows that
\[
 \|R_i\|_F\le\frac83\sqrt m+\frac23\tau m.
\]
Use the pivot supplied by Lemma~\ref{gram} and rearrange to obtain the
second inequality.
\end{proof}

\paragraph{Restriction and repeated blocks.}
The result applies after restricting all permutations to any common
subset of labels, with exactly the same weights and the same
$\kappa$. If two restrictions coincide, they may be left as repeated
blocks; merging them only increases $\kappa$ and strengthens the bound.
This makes greedy packing of disjoint six-label windows possible
whenever a fixed positive fraction of $N$ remains and
$N\ge C/\kappa$.
\end{document}
